% sections/evaluation.tex --- Empirical Evaluation
\section{Empirical Evaluation}
\label{sec:evaluation}

This section reports the headline results of the released
configuration of \bench{}.  Numerical entries marked
``\textit{[pending training run]}'' will be filled in from the
v$1.0$ release training runs; the methodology, baselines, and
evaluation harness are fully specified here.

\subsection{Released configuration}
\label{sec:evaluation:config}

The released v$1.0$ configuration of \bench{} produces a
twelve-month synthetic enterprise dataset under the
manufacturing industry profile (Section~\ref{sec:datasynth}).
Headline volumes are summarised in Table~\ref{tab:eval:volume};
the verbatim YAML is reproduced in
Appendix~\ref{sec:appendix:config}.

\begin{table}[H]
  \centering
  \caption{Released configuration v$1.0$: headline volumes.
    Edges are produced by Method~A on $1$-to-$1$ entries and
    Method~B on $n$-to-$n$ entries; counts are deterministic
    given the configuration seed.}
  \label{tab:eval:volume}
  \begin{tabular}{lr}
    \toprule
    Period             & FY2024 (12 months)\\
    Industry profile   & Manufacturing\\
    Subsidiaries       & $25$\\
    Chart-of-accounts size (mean) & $\sim$$520$ accounts\\
    Journal entries    & \textit{[pending training run]}\\
    Edge count (\texttt{je\_network.csv}) & \textit{[pending training run]}\\
    Anomaly rate (T1 positive)            & $\sim$$3.5\%$\\
    Fraud rate (subset)                   & $\sim$$1.0\%$\\
    \bottomrule
  \end{tabular}
\end{table}

The fraud-rate split between line-level and document-level fraud
follows the calibration formula of \citep{ivertowski2026datasynth}:
\texttt{fraud\_rate}~$= 0.020$, \texttt{document\_fraud\_rate}~$= 0.050$,
\texttt{propagate\_to\_lines}~$= \mathrm{true}$, yielding the
$\sim$$3.5\%$ headline.

\subsection{Baselines}
\label{sec:evaluation:baselines}

We provide a baseline suite that spans the methodological space of
edge-level anomaly detection.  All baselines train from a single
declarative configuration (\texttt{baselines/config.yaml}); the
hyperparameter tables for each are listed in
Appendix~\ref{sec:appendix:hyperparams}.

\begin{itemize}
  \item \textbf{Graph neural networks.}
    GCN~\citep{kipf2017semi},
    GraphSAGE~\citep{hamilton2017inductive},
    GAT~\citep{velickovic2018graph},
    GIN~\citep{xu2019gin},
    RGCN~\citep{schlichtkrull2018rgcn}.  All are trained as
    edge-classification models with two $128$-dimensional
    hidden layers; node features are constructed from
    account-class one-hots, account-degree statistics, and a
    normalised account-frequency embedding.
  \item \textbf{Tabular gradient boosting.}
    XGBoost~\citep{chen2016xgboost} on a feature vector
    composed of the $13$ edge columns plus aggregated
    neighbour statistics (mean / max / std of incident-edge
    amounts and same-document edge counts).
  \item \textbf{Classical anomaly detectors.}
    Local Outlier Factor~\citep{breunig2000lof} and
    Isolation Forest~\citep{liu2008isolation} on the same
    feature vector as XGBoost, used unsupervised and scored
    against the ground-truth labels.
  \item \textbf{Rule-based reference.}  An ISA~240-style rule
    set~\citep{isa240}: round-dollar amounts, weekend or
    off-hours postings, post-close postings, manual postings to
    typically-automated accounts.  This baseline is the
    fairest non-ML reference for ``what an audit programme
    catches today''.
\end{itemize}

\subsection{Hyperparameter protocol}
\label{sec:evaluation:hp}

For each baseline we use the validation slice of the temporal
split to select hyperparameters.  We do not perform a full grid
search; we report the single configuration shipped with the
release plus, for the GNN family, a learning-rate sweep over
$\{10^{-3}, 5\!\times\!10^{-4}, 10^{-4}\}$ and a depth sweep over
$\{2, 3\}$ layers.  The fixed configuration is what the release
trains; the sweeps appear only in the ablation Section
\ref{sec:evaluation:ablations}.

\subsection{Headline results}
\label{sec:evaluation:headline}

Table~\ref{tab:eval:headline} reports primary $\AUCPR$ on T1 for
each split.  Entries marked \textit{[pending]} will be supplied
in the v$1.0$ release.

\begin{table}[H]
  \centering
  \caption{T1 binary anomaly detection: primary $\AUCPR$ by split.
    Higher is better.  $95\%$ bootstrap CIs in parentheses.
    \textit{[All values pending v$1.0$ training run.]}}
  \label{tab:eval:headline}
  \small
  \begin{tabular}{lccc}
    \toprule
    \textbf{Method} & \textbf{Temporal} & \textbf{Entity} &
      \textbf{Typology}\\
    \midrule
    Rule-based reference   & \textit{[pending]} & \textit{[pending]} & \textit{[pending]}\\
    Local Outlier Factor   & \textit{[pending]} & \textit{[pending]} & \textit{[pending]}\\
    Isolation Forest       & \textit{[pending]} & \textit{[pending]} & \textit{[pending]}\\
    XGBoost                & \textit{[pending]} & \textit{[pending]} & \textit{[pending]}\\
    \midrule
    GCN~\citep{kipf2017semi}              & \textit{[pending]} & \textit{[pending]} & \textit{[pending]}\\
    GraphSAGE~\citep{hamilton2017inductive} & \textit{[pending]} & \textit{[pending]} & \textit{[pending]}\\
    GAT~\citep{velickovic2018graph}        & \textit{[pending]} & \textit{[pending]} & \textit{[pending]}\\
    GIN~\citep{xu2019gin}                 & \textit{[pending]} & \textit{[pending]} & \textit{[pending]}\\
    RGCN~\citep{schlichtkrull2018rgcn}    & \textit{[pending]} & \textit{[pending]} & \textit{[pending]}\\
    \bottomrule
  \end{tabular}
\end{table}

We expect, on the basis of preliminary internal runs, that the
ranking on the temporal split groups roughly into three tiers:
relational-aware GNNs (RGCN, GAT) at the top, attention-free GNNs
and XGBoost in a competitive middle band, and the unsupervised
classical detectors and the rule-based reference forming the
floor.  We expect the entity-split numbers to lag the temporal
numbers by a small margin and the typology-split numbers to lag
substantially---this gap is itself one of the headline findings
the benchmark is designed to surface.

\subsection{Per-typology breakdown}
\label{sec:evaluation:typology}

Table~\ref{tab:eval:typology} reports T2 per-typology $\AUCPR$ on
the temporal split.  We anticipate, again from preliminary runs,
that fraud is the most learnable category (it is also the
category with the most distinctive behavioural-bias signature),
followed by error and process anomalies; statistical and
relational categories are harder and benefit most from network
features.

\begin{table}[H]
  \centering
  \caption{T2 multi-class $\AUCPR$ on the temporal split, by
    anomaly category.  \textit{[All values pending v$1.0$ training
    run.]}}
  \label{tab:eval:typology}
  \small
  \begin{tabular}{lccccc}
    \toprule
    \textbf{Method} & \textbf{Fraud} & \textbf{Error} &
      \textbf{Process} & \textbf{Statistical} & \textbf{Relational}\\
    \midrule
    Rule-based reference & \textit{[pending]} & \textit{[pending]} & \textit{[pending]} & \textit{[pending]} & \textit{[pending]}\\
    XGBoost              & \textit{[pending]} & \textit{[pending]} & \textit{[pending]} & \textit{[pending]} & \textit{[pending]}\\
    Best GNN             & \textit{[pending]} & \textit{[pending]} & \textit{[pending]} & \textit{[pending]} & \textit{[pending]}\\
    \bottomrule
  \end{tabular}
\end{table}

\subsection{Ablations}
\label{sec:evaluation:ablations}

The release supports three planned ablation experiments, each of
which a downstream user can re-run from configuration:

\begin{itemize}
  \item \textbf{Network depth.}  Hold the GNN family fixed and
    vary the number of message-passing layers in
    $\{1, 2, 3, 4\}$; report $\AUCPR$ on the temporal split.
  \item \textbf{\texttt{predecessor\_edge\_id} feature.}  Re-train
    the best GNN with and without document-chain features
    derived from \texttt{predecessor\_edge\_id}; report the
    delta $\AUCPR$ as evidence for the value of chain context.
  \item \textbf{Structure-only vs.\ attribute-only.}  Train two
    additional XGBoost variants: one on edge-list structure
    (no amounts, no dates, no flags) and one on edge attributes
    only (no graph features); both are reported alongside the
    full XGBoost as a decomposition of the signal.
\end{itemize}

\subsection{Cross-typology generalisation}
\label{sec:evaluation:cross}

We also expect the typology-hold-out split to be the hardest of
the three.  Initial runs on a near-final configuration showed all
baselines falling within a narrow $\AUCPR$ band on T3---substantially
below their T1 performance on the same data.  This finding, which
we will quantify in the v$1.0$ release, is a candidate driver for
methodological work on \emph{class-conditional} fraud detection
that does not currently dominate the literature.
